\chapter{Discusión}

\section{Interpretación física}
Los patrones observados reflejan competencia no recíproca: $c$ domina y fragmenta el espacio cuando $\mu=0$, mientras que $\mu=1$ introduce una energía libre que ordena dominios y reduce heterogeneidad. El efecto Allee fuerte actúa como barrera crítica que limita la extensión tumoral aun con control moderado.

\section{Comparación con la literatura}
Los resultados son coherentes con modelos de inmunoterapia \cite{delitala_mathematical_2007} y adhesión celular \cite{gatenby_role_2004}, donde la geometría de nulclinas y el espectro lineal definen umbrales de invasión. La magnitud de Re\,$\lambda_{\max}$ en Weak vs Strong subraya la sensibilidad del sistema a la dureza del Allee.

\section{Implicaciones}
\begin{itemize}
  \item No se hallaron equilibrios estables en barridos actuales; el sistema opera cerca de umbrales inestables.
  \item El control variacional ($\mu=1$) organiza patrones y puede reducir la carga tumoral, pero no elimina inestabilidad.
  \item El efecto fuerte sugiere ventanas de intervención temprana: pequeñas perturbaciones pueden impedir despegue si no se supera el umbral.
\end{itemize}

\section{Limitaciones}
Faltan corridas con mallas más finas y rangos extendidos de parámetros para buscar estabilidad; los protocolos $u(x,y,t)$ no se han optimizado y las correlaciones se evaluaron sólo en regímenes cortos ($T\le 2$).




